\section{\secconvcp}
\label{sec_converse_cp}

In this section, we continue our discussion of conditional statements, including the converse and contrapositive of a conditional statement.  The contrapositive provides an alternative proof format for proving that a conditional is true.  We also introduce biconditional (if and only if) statements and the proof format for biconditional statements.  % Su: want to mention any applications here?

%%%%%%%%%%%%%%%%%%%%%%%
\newcommand{\subconv}{The Converse and Biconditional Statements}
\subsection{\subconv}
\label{sub_converse}

The statement $P \to Q$ promises that if $P$ is true, then $Q$ is also true.  It does not say the reverse.  The ``reverse'' might be one of two statements, each of which has a name.

\begin{defn}[Converse and inverse]
\label{defn_converse}
~
    \begin{apart}
        \item For statements $P$ and $Q$, the \vocab{converse} of $P \to Q$ is $Q \to P$.  For example, for the statement 
            \begin{center}
                For all integers $s$, if $s$ is even, then $s^2$ is even.
            \end{center}
        the converse is
            \begin{center}
                For all integers $s$, if $s^2$ is even, then $s$ is even.
            \end{center}
        \item In this context, we can translate the conditional $P \to Q$ as ``$Q$ \vocab{if} $P$'' and we can translate the converse $Q \to P$ as ``$Q$ \vocab{only if} $P$''.
            
        \item For statements $P$ and $Q$, the \vocab{inverse} of $P \to Q$ is $\neg P \to \neg Q$.  For example, for the statement 
            \begin{center}
                For all integers $m$, if $m$ is even, then $m^2$ is even.
            \end{center}
        the inverse is
            \begin{center}
                For all integers $m$, if $m$ is odd, then $m^2$ is odd.
            \end{center}
    \end{apart}
\end{defn}

Let's compare the truth table for a conditional statement to its converse.

\begin{exam}[The converse]
\label{exam_converse}
~
    \begin{apart}
        \item Construct a truth table for a conditional statement and its converse.
        
            \begin{soln}
                We constructed the truth table for the converse in \Cref{exam_truthtable_compound_conditionals}.  We repeat that here alongside the conditional for comparison in \Cref{table_truthtable_converse}. Notice that both a statement and its converse can be true or either can be true with the other false.
            \end{soln}

            \begin{table}[ht]
                \centering
                \begin{tabular}{cc|c|c}
                     $P$ & $Q$ & $P \to Q$ & $Q \to P$ \\ \hline
                     \str{T} & \str{T} & \str{T} &\str{T}  \\
                     \str{T} & \str{F} & \str{F} &\str{T}  \\
                     \str{F} & \str{T} & \str{T} &\str{F}  \\
                     \str{F} & \str{F} & \str{T} &\str{T}  \\
                \end{tabular}
                \caption{Truth table for a conditional and its converse.}
                \label{table_truthtable_converse_again}
            \end{table}
            
            \item Give an example of a true conditional statement that has a false converse.
                \begin{soln}
                    Consider the conditional statement 
                        \begin{center}
                            If $k=3$, then $k^2 = 9$.
                        \end{center}
                    This statement is true because if $k=3$, then $k^2=3^2=9$.
                    The converse is 
                        \begin{center}
                            if $k^2=9$, then $k=3$.
                        \end{center}
                    The converse is false because if $k=-3$, then $k^2=(-3)^2=9$ is true but $k=-3=3$ is false. That is, we have a broken promise (\Cref{thm_broken_promise}).
                \end{soln}
            \item \label{part_example_true_converse} Give an example of a true conditional statement that has a true converse.
                \begin{soln}
                    Consider the conditional statement 
                        \begin{center}
                            If $n$ is odd, then $n+1$ is even.
                        \end{center}
                    The converse is 
                       \begin{center}
                            If $n+1$ is even, then $n$ is odd.
                        \end{center}
                    We proved that both statements are true in 
                    \Cref{sec_even_odd_divides}  \Cref{exer_prove_consecutive_even_odd} and again in \Cref{sec_conditional_statements} \Cref{exer_integer_or_next_is_even}. 
                \end{soln}
    \end{apart}
\end{exam}

When both a conditional statement, $P \to Q$, and its converse, $Q \to P$, are true, we say ``$P$ \vocab{if and only if} $Q$''.  We give a formal definition of the biconditional.

\begin{defn}[Biconditional statements]
\label{defn_biconditional}
~
    \begin{apart}
    \item For statements $P$ and $Q$, the \vocab{biconditional} statement, denoted $P \leftrightarrow Q$, is the following compound statement:
        $$(P \to Q) \land (Q \to P).$$
    For example, we saw in \Cref{exam_converse} (\ref{part_example_true_converse}) that if $P$: $n$ is odd and $Q$: $n+1$ is even, then $P \to Q$ and $Q \to P$ are true, so the biconditional statement $P \leftrightarrow Q$ is also true. 
    
    \item The biconditional $P \leftrightarrow Q$ is often translated as $P$ \vocab{if and only if} $Q$ and may also be abbreviated $P$ iff $Q$.  For example, when $P$: $n$ is odd and $Q$: $n+1$ is even, the biconditional $P \leftrightarrow Q$ says 
        \begin{center}
            For any integer $n$, $n$ is odd if and only if $n+1$ is even.
        \end{center}
    \end{apart} 
\end{defn}

Let's construct a truth table for a biconditional statement.

\begin{exam}[Truth table of biconditional statement]
\label{exam_truthtable_biconditional}~
~
    \begin{apart} 
        \item Construct a truth table for a biconditional statement.
            \begin{soln}
                We construct the truth table in \Cref{table_truthtable_biconditional}.
            \end{soln}

        \begin{table}[hbtp]
            \centering
                \begin{tabular}{cc|ccc}
                     $P$ & $Q$ & $P \to Q$ & $Q \to P$ & $P \leftrightarrow Q$\\ \hline
                     \str{T} & \str{T} & \str{T} &\str{T} &\str{T}   \\
                     \str{T} & \str{F} & \str{F} &\str{T}  &\str{F}  \\
                     \str{F} & \str{T} & \str{T} &\str{F} &\str{F}   \\
                     \str{F} & \str{F} & \str{T} &\str{T}  &\str{T}  \\
                \end{tabular}
            \caption{Truth table for the biconditional $P \leftrightarrow Q$.}
            \label{table_truthtable_biconditional}
        \end{table}

        \item Use this truth table to explain when the biconditional $P \leftrightarrow Q$ is true and when it is false.
            \begin{soln}
                The biconditional $P \leftrightarrow Q$ is true when $P$ and $Q$ are both true (row 1) or $P$ and $Q$ are both false (row 4). That is, $P \leftrightarrow Q$ is true when $P$ and $Q$ have the same truth value meaning $P \equiv Q$. On the other hand, the biconditional $P \leftrightarrow Q$ is false when $P$ is true but $Q$ is false (row 2) or when $P$ is false but $Q$ is true (row 3). That is, $P \leftrightarrow Q$ is false when $P$ and $Q$ have the opposite truth value meaning $Q \equiv \neg P$.
            \end{soln} 
    \end{apart}
\end{exam}

Let's practice working with the converse and biconditional statements.

\begin{act}[The converse and biconditional statements]
\label{act_converse_iff}
~
    \begin{actpart}
        \item State the converse of the following statement: if the graph $G$ has a vertex of degree five, then the graph $G$ has at least six vertices. Is the statement true?  Is the converse true?
        \item \label{part_converse_105} State the converse of the following statement: for any integer $k$, if $k >10$, then $k > 5$. Is the statement true?  Is the converse true?
        \item \label{part_pos_one} State the converse of the following statement:  for any integer $t$, if $t \ge 1$, then $t >0$.  Is the statement true? Is the converse true?  Hint: $t$ is an integer.
        \item \label{part_iff_nsqmult4} True or false?  For any integer $m$, $4 \mid m^2$ if and only if $4 \mid m$.
        \item \label{part_iff_3neven} True or false?  For any integer $n$, $n$ is even if and only if $3n$ is even.
    \end{actpart}
\end{act}

%%%%%%%%%%%%%%%%%%%%%%%
\newcommand{\subpffiff}{Proof Format: $\leftrightarrow$}
\subsection{\subpffiff}
\label{sub_pff_iff}

Since the biconditional is defined as $$P \leftrightarrow Q \equiv (P \to Q) \land (Q \to P),$$ we can combine the proof formats for ``and'' (\Cref{pff_and}) with the proof format for a conditional statement (\Cref{pff_if_then_direct}) to get a new proof format to prove a biconditional statement.  Normally in ``and'' proofs we begin the two parts with the words ``First'' and ``Second'' (or ``First'' and ``Next''). In the proof of a biconditional statement, the second part proves the converse, so we begin with the word ``Conversely'' instead.

\begin{pff}[Proof of $\leftrightarrow$]
\label{pff_iff} 
~
    Prove: $P \leftrightarrow Q$ 
    \begin{proof}
        First, assume $P$. \\
        (Explain why $Q$ is true.) \\
        Thus, $Q$. \\
        Conversely, assume $Q$. \\
        (Explain why $P$ is true.) \\
        Thus, $P$.
    \end{proof}    
\end{pff}

As usual, when we introduce a new proof format, we provide an example for you to copy and edit.

\begin{act}[Biconditional proof]
\label{act_iff_proof}
~    
    \begin{actpart}
        \item Copy the following proof.

        Prove for any integer $a$ that $a \fn{ mod } 3 = 1$ if and only if $3 \mid (a-1)$.

        \emph{Proof.}  Let $a$ be an integer.

        First, assume that $a \fn{ mod } 3 = 1$. By definition of \fn{mod}, it follows that $a = 3q +1$ for some integer $q$.  Subtracting one from each side of this equation we get $a-1 = 3q$.  By definition of divides, it follows that $3 \mid (a-1)$.

        Conversely, assume that $3 \mid (a-1)$. By definition of divides, we can write $a-1=3k$ for some integer $k$.  Adding one to each side of this equation we get $a = 3k+1$.  By definition of \fn{mod}, it follows that $a \fn{ mod } 3 = 1$.

        \item Edit the proof to prove for any integers $a$ and $n$ that $a \fn{ mod } n = 1$ if and only if $n \mid (a-1)$.
    \end{actpart}
\end{act}

% SU:  you could do the condensed format, but I think it's okay to skip.
% \begin{align*}
% a \fn{ mod }3 = 1 & \leftrightarrow a = 3q+1 \text{ for some integer }q \\
% & \leftrightarrow a-1 = 3q \text{ for some integer }q \\  
% & \leftrightarrow 3 \mid (a-1)
% \end{align*}


We can now prove the connection between divides and fractions.

\begin{exam} [Fractions and Divisibility]
\label{exam_fractions_divides}
    Let $d$ and $n$ be integers with $d \neq 0$. Prove that $d \mid n$ exactly when $\frac{n}{d}$ is an integer.
    
        \emph{Proof.}  Let $d$ and $n$ be integers with $d \neq 0$.
        
        \quad First, assume $d \mid n$. By definition of divides, $n=dk$ for some integer $k$.  By \Cref{thm_equal_fractions}, it follows that
            $$\frac{n}{d} = \frac{dk}{d(1)} = \frac{k}{1} = k$$
        is an integer.
            
        \quad Conversely, assume $\frac{n}{d}$ is an integer.  Say, 
            $$\frac{n}{d} = k = \frac{k}{1}.$$  
        By \Cref{thm_equal_fractions}, it follows that $n(1) = k(d)$ or, equivalently, $n=dk$.  By the definition of divides, $d \mid n$. \hfill $\Box$          
\end{exam}

%%%%%%%%%%%%%%%%%%%%%%%
\newcommand{\subcp}{The Contrapositive}
\subsection{\subcp}
\label{sub_contrapositive}

While a conditional statement and its converse are not logically equivalent, there is a related conditional that is logically equivalent to a conditional.

\begin{defn}[Contrapositive]
\label{defn_contrapositive} 
    For statements $P$ and $Q$, the \vocab{contrapositive} of the conditional statement $P \to Q$ is the conditional statement $\neg Q \to \neg P$. For example, consider the true statement 
        \begin{center}
            If $k=3$, then $k^2=9$.
        \end{center}
    then the contrapositive is
        \begin{center}
            If $k^2 \neq 9$, then $k \neq 3$.
        \end{center}
    which is also true.
\end{defn}

Let's explore the logical equivalence of a conditional and its contrapositive.

\begin{exam}[Equivalence of statement and its contrapositive]
\label{exam_equiv_contrapositive}
~
    \begin{apart}
        \item Construct a truth table of conditional and its contrapositive.  What do you notice?
            \begin{soln}
                The truth table is shown in \Cref{table_truthtable_contrapositive}.  Notice that the columns for the conditional and its contrapositive are identical.  That means that a conditional and its contrapositive are logically equivalent.
            \end{soln}
            \begin{table}[hbtp]
                \centering
                \begin{tabular}{cc|c|ccc}
                $P$ & $Q$ & $P \to Q$ & $\neg Q$ & $\neg P$ & $\neg Q \to \neg P$   \\ \hline  
                \str{T} & \str{T} & \str{T} & \str{F} & \str{F} & \str{T}  \\
                \str{T} & \str{F} & \str{F} & \str{T} & \str{F} & \str{F}  \\
                \str{F} & \str{T} & \str{T} & \str{F} & \str{T} & \str{T} \\
                \str{F} & \str{F} & \str{T} & \str{T}& \str{T} & \str{T}  \\ 
            \end{tabular}
                \caption{Truth table for the contrapositive}
                \label{table_truthtable_contrapositive}
            \end{table}
        \item Explain in words why the contrapositive is logically equivalent to the original conditional.
            \begin{soln}
                Think about what makes the conditional $P \to Q$ is true.  We are promising that if $P$ is true, then $Q$ is true.  What happens if $\neg Q$ is true? That means $Q$ is false and so $P$ cannot be true (or we would break our promise) and so $P$ must be false which means $\neg P$ is true.  That is, if $\neg Q$, then $\neg P$, which is the contrapositive.
            \end{soln}
    \end{apart}   
\end{exam}

We add the contrapositive to our list of logical equivalences.

\begin{thm}[Contraposition equivalence]
\label{thm_cp}
    For statements $P$ and $Q$, 
        $$(P \to Q) \equiv (\neg Q \to \neg P)$$
\end{thm}

Let's practice working with the contrapositive.

\begin{act}[Contrapositive]
\label{act_cp}
~  
    \begin{actpart}    
        \item Write the contrapositive of the true statement: For all integers $n$, if $n>10$, then $n>5$.
        \item Write the contrapositive of the true statement: For all graphs $G$, if $G$ has at most five vertices, then $G$ has at most 10 edges. Hint: To negate ``at more'' use ``more than''.
        \item \label{part_cp_nsqodd_nodd} Write the contrapositive of the true statement: For all integers $n$, if $n^2$ is odd, then $n$ is odd.
        \item \label{part_cp_easier} Which statement do you think would be easier to prove:  the statement from (\ref{part_cp_nsqodd_nodd}) or its contrapositive? Explain.        
        \end{actpart}
\end{act}

%%%%%%%%%%%%%%%%%%%%%%%
\newcommand{\subpffcp}{Proof Format: $\to$ Using the Contrapositive}
\subsection{\subpffcp}
\label{sub_pff_cp}

Since a conditional and its contrapositive are logically equivalent, we have a new way to prove $P \to Q$: we can prove its contrapositive $\neg Q \to \neg P$ instead.  

\begin{pff}[Contrapositive proof of $\to$]
\label{pff_if_then_cp}
    Prove: $P \to Q$
    \begin{proof}
    We will prove the contrapositive, namely (fill in $\neg Q \to \neg P$).  
    
    Assume $\neg Q$ is true. 
    
    (Explain why it follows that $\neg P$ is true.)

    Thus $\neg P$ is true.
    \end{proof}
\end{pff}

Let's do an example with this new proof format.

\begin{act}[Contrapositive proof if $n^2$ is odd, then $n$ is also odd]
\label{act_direct_proof_ifnsqoddthennodd}
~
    \begin{actpart}
        \item \label{part_ifnsqoddthennodd} Perhaps you noticed in \Cref{act_cp} (\ref{part_cp_easier}) that the contrapositive would be easier to prove.  Copy the following proof.

        Prove for any integer $n$, if $n^2$ is odd then $n$ is odd.
        \begin{proof}
            Let $n$ be an integer.  We will prove the contrapositive, namely that if $n$ is even, then $n^2$ is even.         
            
            Assume that $n$ is even.  By the definition of even, we can write $n=2k$ for some integer $k$.  Then $$n^2=(2k)^2=(2k)(2k) = (2)(2)kk = 4k^2 = 2(2k^2).$$
            Since $2k^2$ is an integer, it follows from the definition of even that $n^2$ is even.
        \end{proof}
        \item Edit the proof to prove for any integer $n$ if $5n$ is odd, then $n$ is odd.  
    \end{actpart}
\end{act}

When a proof by contradiction does not use the supposed hypothesis until the ending contradiction, we can often rewrite the proof using the contrapositive, avoiding the proof by contradiction.

\begin{exam}[Rewriting a proof by contradiction using the contrapositive] 
\label{exam_product_even_implies_even_factor}
~
    \begin{apart}
        
        \item \label{part_contradiction_product_even_implies_even_factor} Read the following proof by contradiction.
        
            Prove that for any integers $a$ and $b$ if $ab$ is even, then $a$ is even or $b$ is even. 
            
            \begin{proof}
                Suppose not.  Then for some integers $a$ and $b$ we would have that $ab$ is even but neither $a$ nor $b$ is even.  
                
                First, since neither $a$ nor $b$ is even, that means $a$ and $b$ are odd.  By the definition of odd, we can write $a = 2k_1+1$ and $b = 2k_2+1$ for some integers $k_1$ and $k_2$.  Multiplying $a$ and $b$ we get $$ ab=(2k_1+1)(2k_2+1)=4k_1k_2+2k_1+2k_2+1 = 2(2k_1k_2 + k_1 + k_2) + 1$$ and so $ab$ is odd.

                On the other hand, we said $ab$ was even.  This is a contradiction.
            \end{proof} 

            What are the contradictory statements $Q$ and $\neg Q$?
                \begin{soln}
                    The contradictory statements are $Q: ab$ is odd and $\neg Q: ab$ is even.
                \end{soln}
            
        \item State the contrapositive of the statement ``For any integers $a$ and $b$ if $ab$ is even, then $a$ is even or $b$ is even.'' 
            \begin{soln}
                By DeMorgan's Laws (\Cref{thm_demorgans}), the negation of the conclusion
                    \begin{center}
                        $a$ is even or $b$ is even.
                    \end{center} 
                is the statement
                    \begin{center}
                        $a$ and $b$ are odd.
                    \end{center}
                The contrapositive is 
                    \begin{center}
                        For any integers $a$ and $b$, if $a$ and $b$ are odd, then $ab$ is odd.
                    \end{center}
                Notice that we keep the universal quantifier.  
            \end{soln}
        
        \item Adapt the proof from (\ref{part_contradiction_product_even_implies_even_factor}) to prove for any integers $a$ and $b$ if $ab$ is even, then either $a$ is even or $b$ is even using the contrapositive (\Cref{pff_if_then_cp}).  Hint: you should not have to change much of the proof.
            \begin{soln}
                Prove: for any integers $a$ and $b$ if $ab$ is even, then either $a$ is even or $b$ is even. 
            
                \emph{Proof.}
                We will prove the contrapositive, namely for any integers $a$ and $b$, if $a$ and $b$ are odd, then $ab$ is odd.    

                Let $a$ and $b$ be odd integers.  By the definition of odd, we can write $a = 2k_1+1$ and $b = 2k_2+1$ for some integers $k_1$ and $k_2$.  Multiplying $a$ and $b$ we get $$ ab=(2k_1+1)(2k_2+1)=4k_1k_2+2k_1+2k_2+1 = 2(2k_1k_2 + k_1 + k_2) + 1$$ and so $ab$ is odd.
            
            \end{soln}
    \end{apart}
\end{exam}

%%%%%%%%%%%%%%%%%%%%%%%

\subsection{Exercises}

%%%%%%%%%%%%%%%%%%%%%%%%

\subsubsection*{Exercises for \Cref{sub_converse} \subconv}

\begin{exer}[\exerP] % write converse true/false even
\label{exer_converse_TorF_lessthan}
    Consider the statement: for any real number $x$, if $x^2 \le 100$, then $x \le 10$.
        \begin{exerpart}
            \item Is the statement true or false?
            \item State the converse.
            \item Give a counterexample to show that the converse is false. 
        \end{exerpart}
\end{exer}

\begin{exer}[\exerP] % write converse true/false divides
\label{exer_converse_TorF_divides}
    Consider the statement: for all integers $n$, if $6 \mid n$, then $3 \mid n$.
        \begin{exerpart}
            \item Is the statement true or false?
            \item State the converse.
            \item Give a counterexample to show that the converse is false.
        \end{exerpart}
\end{exer}

\begin{exer}[\exerP] % write converse true/false mod
\label{exer_converse_TorF_mod}
    Consider the statement: for any integers $a$ and $b$, if $a \fn{ mod } 5 = b \fn{ mod } 5$, then $a^2 \fn{ mod } 5 = b^2 \fn{ mod } 5$
        \begin{exerpart}
            \item Is the statement true or false?
            \item State the converse.
            \item Give a counterexample to show that the converse is false.
        \end{exerpart}
\end{exer}

\begin{exer}[\exerP] % show converse and inverse are logically equivalent.
\label{exer_converse_inverse_equiv}
    Use a truth table to show that the converse $Q \to P$ and the inverse $\neg P \to \neg Q$ are logically equivalent.  Include intermediate columns for $\neg P$ and $\neg Q$.  This result helps explain why we rarely discuss the inverse.
\end{exer}

\begin{exer}[\exerP] % decide biconditional t/f
\label{exer_TorF_iff}   
    Decide if each biconditional statement (from \Cref{act_converse_iff}) is true or false.  Explain briefly.
        \begin{exerpart}
            \item For any integers $k$, $k>10$ if and only if $k>5$.
            \item For any integer $n$, $n$ is even if and only if $3n$ is even.
            \item For any integer $m$, $4 \mid m^2$ if and only if $4 \mid m$.
        \end{exerpart}
\end{exer}

\begin{exer}[\exerU] % negating a biconditional
\label{exer_neg_biconditional}   
~
    \begin{exerpart}
        \item Construct a truth table for $\neg(P \leftrightarrow Q)$.  Hint: Copy the truth table from \Cref{exam_truthtable_biconditional} and add one additional column.      
        \item The negation $\neg(P \leftrightarrow Q)$ is logically equivalent to a simpler statement.  What is it?  Hint: \Cref{table_truthtable_xor}.
    \end{exerpart}
\end{exer}

\begin{exer}[\exerR] % Converse and Biconditionals
\label{exer_dyk_converse_iff}
    Do you know \ldots
        \begin{exerpart}
            \item What the converse of $P \to Q$ is?
            \item How to find examples to show that $P \to Q$ and its converse are not equivalent?
            \item How the biconditional statement $P \leftrightarrow Q$ is defined?
            \item When a biconditional statement is true and when it is false?
        \end{exerpart}
\end{exer}

%%%%%%%%%%%%%%%%%%%%%%%%

\subsubsection*{Exercises for \Cref{sub_pff_iff} \subpffiff}

\begin{exer}[\exerP] % even
\label{exer_proof_iff_even}
    Use \Cref{pff_iff} to prove that $n$ is even if and only if $n+3$ is odd, for any integer $n$.
\end{exer}
 
\begin{exer}[\exerP] %solution equation
\label{exer_proof_iff_solution}
    Use \Cref{pff_iff} to prove that $2n-7=5$ if and only if $n =6$ for any integer $n$.
\end{exer}

\begin{exer}[\exerR] % Proof format biconditional
\label{exer_dyk_pff_iff}
    Do you know \ldots
        \begin{exerpart}
            \item How to format the proof of a biconditional statement?
            \item Why the second part of the proof of a biconditional statement begins with the word ``Conversely''?
        \end{exerpart}
\end{exer}

% Many examples use the contrapositive to prove the converse so really have to wait. Potentially these could be two sections, but I like it like this.

%%%%%%%%%%%%%%%%%%%%%%%%

\subsubsection*{Exercises for \Cref{sub_contrapositive} \subcp}

\begin{exer}[\exerP] %  write contrapositive]
\label{exer_cp_assorted}
Write the contrapositive of each statement.
    \begin{exerpart}
        \item For any real number $x$, if $3x \le 6$, then $x \le 2$.
        \item For any integer $s$, if $s^2$ is even, then $s$ is even.
        \item For all graphs $G$, if $G$ has at least 5 vertices, then $G$ must have at least 5 edges.
    \end{exerpart}
\end{exer}

\begin{exer}[\exerU] %  subgraph of acyclic graph]
\label{exer_subgraph_acyclic}
    Let $H$ be a subgraph of a graph $G$. Consider the statement
        \begin{center} If $G$ is acyclic, then $H$ is acyclic. \end{center}
        \begin{exerpart}
            \item State the contrapositive of this statement and explain why it's true.
            \item Explain why it follows that this statement is true.  %Note:  we used this statement in the proof in \Cref{exam_math_ind_tree_size} when we wrote ``Since $T$ was a tree, it does not contain any cycles so neither does $T'$.''
        \end{exerpart}
\end{exer}

\begin{exer}[\exerU] %  write contrapositives involving connectives]
\label{exer_cp_irrat}
Write the contrapositive of each statement.
    \begin{exerpart}
        \item For all real numbers $x$ and $y$, if $x$ is rational and $y$ is irrational, then $x+y$ is irrational.  Note: from \Cref{defn_sets_of_real_numbers} rational and irrational are negations.
        \item For all real numbers $x$ and $y$, if $x+y$ is irrational, then either $x$ is irrational or $y$ is irrational.
        \item For all real numbers $x$ and $y$, if $x+y$ is rational, then $x$ is rational and $y$ is rational.       
    \end{exerpart}
\end{exer}

\begin{exer}[\exerR] % Contrapositive
\label{exer_dyk_cp}
    Do you know \ldots
        \begin{exerpart}
            \item What the contrapositive of $P \to Q$ is?
            \item Why a conditional statement and its contrapositive are equivalent?
            \item How to use a truth table to verify that a conditional statement and its contrapositive are logically equivalent?
        \end{exerpart}
\end{exer}

%%%%%%%%%%%%%%%%%%%%%%%%

\subsubsection*{Exercises for \Cref{sub_pff_cp} \subpffcp}

\begin{exer}[\exerP] % prove if sq even, then int even.
\label{exer_proof_sq_even_implies_even}
    Use the contrapositive (\Cref{pff_if_then_cp}) to prove that for any integer $s$, if $s^2$ is even then $s$ is even.  
\end{exer}

\begin{exer}[\exerP] % prove 3m odd, then m odd.
\label{exer_proof_3m_odd}
    Use the contrapositive (\Cref{pff_if_then_cp}) to prove for any integer $m$, if $3m$ is odd, then $m$ is odd.  (We proved the converse in \Cref{act_pffifthen} (\ref{part_if_then_direct}).)
\end{exer}
    
\begin{exer}[\exerU] % prove \Cref{thm_show_is_prime} (\ref{part_only_check_prime_divisors})]
\label{exer_prove_only_check_prime_divisors}
    In this exercise we prove \Cref{thm_show_is_prime} (\ref{part_only_check_prime_divisors}).

     Use the contrapositive (\Cref{pff_if_then_cp}) to prove that for any integers $n$, $p$, and $m$, if $p \nmid n$, then $mp \nmid n$. Note that this statement can be rephrased as if $p$ is not a divisor of $n$, then no multiple of $p$ is a divisor of $n$.  
\end{exer}

\begin{exer}[\exerU] %  prove \Cref{thm_show_is_prime} (\ref{part_only_check_to_sqrtn})]
\label{exer_prove_only_check_to_sqrtn}
    In this exercise we prove \Cref{thm_show_is_prime} (\ref{part_only_check_to_sqrtn}).

    Prove for any integers $a$, $b$, and $n$, if $ab=n$, then $a \le \sqrt{n}$ or $b \le \sqrt{n}$.  Hint: Prove the contrapositive (\Cref{pff_if_then_cp}).
\end{exer}

\begin{exer}[\exerR] % Proof conditional using contrapositive
\label{exer_dyk_pff_cp}
    Do you know \ldots
        \begin{exerpart}
            \item How to prove a conditional statement by proving the contrapositive instead?
            \item When we want to prove the contrapositive instead of the original conditional statement?
            \item What we can write in a proof using the contrapositive so that the reader knows we are using the contrapositive?
        \end{exerpart}
\end{exer}

\begin{exer}[\exerE] % proof biconditional using the converse for part
\label{exer_proof_iff_using_cp_evens}
    Use \Cref{pff_iff} to prove that $m$ is odd if and only if $m^3$ is odd.  Hint: Use the contrapositive (\Cref{pff_if_then_cp}) for the converse (second part).  
\end{exer}

\begin{exer}[\exerE] % $\sqrt{2}$ is irrational 
\label{part_prove_sqrtp_irrat} 
    Copy the following proof by contradiction that $\sqrt(2)$ is irrational, being sure to figure out why each statement is true.

        \begin{proof}
            Suppose not.  Then $\sqrt{2}$ is rational.  By the definition of rational (\Cref{defn_sets_of_real_numbers}), we can write $$\sqrt{2} = \frac{s}{t}$$ for some positive integers $s$ and $t$ with $t \neq 0$. By \Cref{thm_equal_fractions}, we may assume that we have canceled any common factors so that $s$ and $t$ are coprime.
            
            On the other hand, multiply each side of this equation by $t$ to get $$s = \sqrt{2}t.$$  Square each side of this equation to get $$s^2=2t^2.$$ Thus, $s^2$ is even.  By \Cref{exer_proof_sq_even_implies_even}, it follows that $s$ is even.  By definition of even,  we can write $s = 2k$ for some integer $k$.  Substituting in we get $$(2k)^2 = 2t^2.$$
            Simplifying we get $t^2=2k^2$. Thus, $t^2$ is also even. By \Cref{exer_proof_sq_even_implies_even}, it follows that $t$ is also even.  Since $s$ and $t$ are even, it follows that two is a common divisor of $s$ and $t$. Thus, $s$ and $t$ are not coprime.
            
            This is a contradiction.  Thus $\sqrt{2}$ is irrational.
        \end{proof}
\end{exer}

%%%%%%%%%%%%%%%%%%%%%%%%%%%
\subsection{\answers}
%%%%%%%%%%%%%%%%%%%%%%%%%%%

%%%%%%%%%%%%%%%%%%
\subsubsection{\answers for \Cref{sub_converse} \subconv}

\subsubsection*{\Cref{exer_converse_TorF_lessthan} \answers}
\begin{apart}
    \item True. Note: In fact, $-10 \le x \le 10$.
    \item Hint:  The converse should still include $\le$.
    \item Hint: Broken promise.
\end{apart}

\subsubsection*{\Cref{exer_converse_TorF_mod} \answers}
\begin{apart}
    \item True
    \item Hint: The converse should also involve $=$.
    \item Hint: Broken promise.
\end{apart}

\subsubsection*{\Cref{exer_TorF_iff} \answers}
\begin{apart}
    \item Hint: One direction is true, but the other direction is false.  Find a counterexample.
    \item Hint: It is true.  Check both directions.
    \item Hint: One direction is true, but the other direction is false.  Find a counterexample.
\end{apart}

%%%%%%%%%%%%%%%%%%
\subsubsection{\answers for \Cref{sub_pff_iff} \subpffiff}

\subsubsection*{\Cref{exer_proof_iff_even} \answers}
    Use \Cref{pff_iff} to prove that $n$ is even if and only if $n+3$ is odd, for any integer $n$.

\subsubsection*{\Cref{exer_proof_iff_solution} \answers}

Hint: Here is an outline of the proof.
     \begin{quote}
        \emph{Proof.}  Let $n$ be an integer.  
        
        First, assume $n$ is even.  By definition of even we can write $n=$ \ldots.  Then $$n+3 = \ldots = 2k'+1$$
        where $k' = \ldots$ is an integer.  Thus $n+3$ is odd.

        Next, assume $n+3$ is odd.  By definition of odd we can write $n+3=$ \ldots.  Then $$n = \ldots = 2k'$$
        where $k' = \ldots$ is an integer.  Thus $n$ is even.    
     \end{quote}

Hint: Here is an outline of the proof.
     \begin{quote}
        \emph{Proof.}  Let $n$ be an integer.  
        
        First, assume $2n-7=5$.  Then, \ldots.  Thus $n=6$.

        Next, assume $n=6$.  Then, $2n-7 = 2(6)-7 = \ldots$.   
     \end{quote}

%%%%%%%%%%%%%%%%%%
\subsubsection{\answers for \Cref{sub_contrapositive} \subcp}

\subsubsection*{\Cref{exer_cp_assorted} \answers}
\begin{apart}
    \item For any real number $x$, if $x > 2$, then $3x > 6$.
    \item Hint: The negation of even is odd.
    \item Hint: The negation of ``at least'' is ``fewer than''.
\end{apart}

\subsubsection*{\Cref{exer_subgraph_acyclic} \answers}
\begin{apart}
    \item Hint: The negation of ``acyclic'' is ``contains a cycle.''
    \item If $H$ contains a subgraph $C$ that is a cycle, then $C$ is a subgraph of $G$ and so $G$ contains a cycle.
\end{apart}

\subsubsection*{\Cref{exer_cp_irrat} \answers}
Hints: First, ``rational'' and ``irrational'' are negations. Also, use DeMorgan's laws any time you need to negate a statement involving ``and'' or ``or''.
\begin{apart}
    \item Hint: The contrapositive begins ``For all real numbers $x$ and $y$, if $x+y$ is rational, then \ldots''
    \item Hint: Your answer should include the word ``and''.  It should not include the word ``irrational''.
    \item Hint: Your answer should not include the word ``rational''.
\end{apart}

%%%%%%%%%%%%%%%%%%
\subsubsection{\answers for \Cref{sub_pff_cp} \subpffcp}

\subsubsection*{\Cref{exer_proof_3m_odd} \answers}
Here is an outline of the proof.
    \begin{quote}
        \emph{Proof.} Let $m$ be an integer.  We will prove the contrapositive, namely \ldots.

        Assume, \ldots.  Then, \ldots.  Thus, \ldots. \hfill $\Box$      
    \end{quote}

\subsubsection*{\Cref{exer_prove_only_check_prime_divisors} \answers}
Here is an outline of the proof.
    \begin{quote}
        \emph{Proof.} Let $n$, $p$, and $m$ be integers.  We will prove the contrapositive, namely if $mp \mid n$, then $p \mid n$.

        Assume, $mp \mid n$.  Then, by definition of divides we can write $n= \ldots = pk$ where $k=\ldots$ is an integer.  Thus, $p \mid n$. \hfill $\Box$      
    \end{quote}

\subsubsection*{\Cref{exer_prove_only_check_to_sqrtn} \answers}
Hint: Use DeMorgan's laws.  Recall that $\sqrt{n}\sqrt{n} = n$.

\subsubsection*{\Cref{exer_proof_iff_using_cp_evens} \answers}
Hint: First, prove if $m$ is odd, then $m^3$ is odd. Conversely, prove the contrapositive of the converse, namely if $m$ is even, then $m^3$ is even.  Another hint: $(2k+1)^3 = (2k+1)(2k+1)(2k+1) = (4k^2+4k+1)(2k+1)=\ldots$












